\section{通用机器人路线：从 Automation 到 Superhuman}

本章总结谭杰对机器人发展阶段的划分。他提到五个阶段：Automation、Teleoperation、Narrow Generalist、Home Generalist、Superhuman Capability。这个路线图有助于避免把所有机器人能力混成一个问题。

\begin{figure}[H]
\centering
\includegraphics[width=0.96\textwidth]{robot-generalist-stages.png}
\caption{通用机器人阶段：从固定自动化到超人能力的长期路线。自制概念图，依据 01:17:35--02:06:16 对谈内容整理。}
\end{figure}

\begin{knowledgebox}{读图：每一阶段的“通用性”都不同}
Automation 是固定规则，Teleoperation 是硬件可用但人远程控制，Narrow Generalist 是窄域智能泛化，Home Generalist 是家庭多任务通用，Superhuman 是在某些能力上超越人类。不要把窄域泛化误读成家庭通用。
\end{knowledgebox}

\subsection{Specialist 与 Generalist}

早期落地可能仍在制造、物流、超市、折衣服、餐巾等垂直场景。垂直 specialist 更容易找到需求、ROI 和可控环境。但如果真正 generalist 成型，它能做 specialist 的任务，还能做更多任务，许多 specialist 会被压缩。问题在于 generalist 需要更久时间和更多数据。

\begin{warningbox}{不要用最终形态否定中间商业化}
通用人形机器人可能是长期目标，但短期垂直场景仍可能创造价值。专业化路线和通用路线不是非此即彼，二者可能并行很多年。
\end{warningbox}

\subsection{安全问题不是儿戏}

谭杰明确说需要关注 AI safety 和 robot safety。当 AI 或机器人能自我迭代时，人类会面临生存问题。Google DeepMind 有 Responsibility and Safety Council，会审查模型和机器人对社会的影响以及安全后果。最坏情况下，如果机器人能力超过安全理解，就应该停下能力扩张，让安全研究追上。

\begin{importantbox}{机器人安全原则}
能力和安全必须齐头并进。每个发展阶段都要做相应安全研究；如果能力进展超过安全理解，应该暂停能力扩张，让 safety catch up。
\end{importantbox}

\subsection{为什么机器人容易被高估}

本节补上谭杰对行业情绪的提醒。机器人最容易被高估，是因为公众看到的常常是最好的 demo：团队可能拍了十遍，选出最好的一遍发布；视频里没有呈现失败次数、环境假设、人工重置、任务边界和安全接管。观众容易把“最优样本”误认为“稳定能力”，从而以为明年就能买到家用人形机器人。

\begin{warningbox}{Demo 视频的证据边界}
一个机器人视频能证明“这件事在某些条件下发生过”，但不能证明它具备可部署能力。要判断落地，需要看重复成功率、失败恢复、环境变化、任务多样性、成本、维护和安全边界。
\end{warningbox}

谭杰的态度是兴奋但冷静：机器人确实在加速，落地场景也开始出现；但“能干活的 humanoid 机器人”现在仍是一片荒漠。真正应该避免的是同时犯两个错误：高估短期能力，低估长期影响。

\begin{importantbox}{正确的时间感}
短期看，机器人离家庭通用还远；中期看，垂直场景可能先落地；长期看，一旦 generalist 真正形成，specialist 的生存空间会被压缩。
\end{importantbox}

\subsection{本章小结}

机器人路线应该分阶段看。短期是窄域泛化和垂直落地，中期是家庭 generalist，长期才是超人能力。安全问题必须从早期就进入研究流程。

